%%%%%%%%%%%%%%%%%%%%%%%%%%%%%%%%%%%%%%%%%
% The Legrand Orange Book
% LaTeX Template
% Version 2.4 (26/09/2018)
%
% This template was downloaded from:
% http://www.LaTeXTemplates.com
%
% Original author:
% Mathias Legrand (legrand.mathias@gmail.com) with modifications by:
% Vel (vel@latextemplates.com)
%
% License:
% CC BY-NC-SA 3.0 (http://creativecommons.org/licenses/by-nc-sa/3.0/)
%
% Compiling this template:
% This template uses biber for its bibliography and makeindex for its index.
% When you first open the template, compile it from the command line with the 
% commands below to make sure your LaTeX distribution is configured correctly:
%
% 1) pdflatex ponteiros_em_c
% 2) makeindex main.idx -s StyleInd.ist
% 3) biber main
% 4) pdflatex main x 2
%
% After this, when you wish to update the bibliography/index use the appropriate
% command above and make sure to compile with pdflatex several times 
% afterwards to propagate your changes to the document.
%
% This template also uses a number of packages which may need to be
% updated to the newest versions for the template to compile. It is strongly
% recommended you update your LaTeX distribution if you have any
% compilation errors.
%
% Important note:
% Chapter heading images should have a 2:1 width:height ratio,
% e.g. 920px width and 460px height.
%
%%%%%%%%%%%%%%%%%%%%%%%%%%%%%%%%%%%%%%%%%

%----------------------------------------------------------------------------------------
%	PACKAGES AND OTHER DOCUMENT CONFIGURATIONS
%----------------------------------------------------------------------------------------

\documentclass[11pt,fleqn]{book} % Default font size and left-justified equations

%%%%%%%%%%%%%%%%%%%%%%%%%%%%%%%%%%%%%%%%%
% The Legrand Orange Book
% Structural Definitions File
% Version 2.1 (26/09/2018)
%
% Original author:
% Mathias Legrand (legrand.mathias@gmail.com) with modifications by:
% Vel (vel@latextemplates.com)
% 
% This file was downloaded from:
% http://www.LaTeXTemplates.com
%
% License:
% CC BY-NC-SA 3.0 (http://creativecommons.org/licenses/by-nc-sa/3.0/)
%
%%%%%%%%%%%%%%%%%%%%%%%%%%%%%%%%%%%%%%%%%

%----------------------------------------------------------------------------------------
%	VARIOUS REQUIRED PACKAGES AND CONFIGURATIONS
%----------------------------------------------------------------------------------------

\usepackage{graphicx} % Required for including pictures
\graphicspath{{Pictures/}} % Specifies the directory where pictures are stored

\usepackage{lipsum} % Inserts dummy text

\usepackage{tikz} % Required for drawing custom shapes

\usepackage[portuguese]{babel} % Portuguese language/hyphenation
\usepackage[utf8]{inputenc} % Required for including letters with accents
\usepackage[T1]{fontenc} % Use 8-bit encoding that has 256 glyphs

\usepackage{enumitem} % Customize lists
\setlist{nolistsep} % Reduce spacing between bullet points and numbered lists

\usepackage{booktabs} % Required for nicer horizontal rules in tables

\usepackage{xcolor} % Required for specifying colors by name
\definecolor{ocre}{RGB}{243,102,25} % Define the orange color used for highlighting throughout the book
\definecolor{azul}{RGB}{13,205,224} % minhas cores
\definecolor{cinza}{RGB}{245,245,245}

\usepackage{listings}
%\usepackage{listingsutf8}
\lstset{language=C, 
	keywordstyle=\color{blue},
	basicstyle=\small\ttfamily,,
	showstringspaces=false,
	xleftmargin=0pt,
	extendedchars=true,
	upquote=true,
	inputencoding=utf8,
	frame=single,
	framerule=0pt,
	breaklines=true,
	emph={True,False},
	texcl=true,
	literate={á}{{\'a}}1 {à}{{\`a}}1 {ã}{{\~a}}1 {â}{{\^a}}1 {é}{{\'e}}1 {í}{{\'i}}1 {ó}{{\'o}}1 {ú}{{\'u}}1 
	{ê}{{\^e}}1  {ô}{{\^o}}1 {õ}{{\~o}}1 {ç}{{\c{c}}}1 {º}{${^\circ}$}1, 
	commentstyle=\color{red},
	tabsize=2,
	stringstyle=\color{red},
	columns = fullflexible, keepspaces=true,
	breakatwhitespace=true,
	backgroundcolor=\color{cinza}
} 

%----------------------------------------------------------------------------------------
%	MARGINS
%----------------------------------------------------------------------------------------

\usepackage{geometry} % Required for adjusting page dimensions and margins

\geometry{
	paper=a4paper, % Paper size, change to letterpaper for US letter size
	top=3cm, % Top margin
	bottom=3cm, % Bottom margin
	left=3cm, % Left margin
	right=3cm, % Right margin
	headheight=14pt, % Header height
	footskip=1.4cm, % Space from the bottom margin to the baseline of the footer
	headsep=10pt, % Space from the top margin to the baseline of the header
	%showframe, % Uncomment to show how the type block is set on the page
}

%----------------------------------------------------------------------------------------
%	FONTS
%----------------------------------------------------------------------------------------

\usepackage{avant} % Use the Avantgarde font for headings
%\usepackage{times} % Use the Times font for headings
\usepackage{mathptmx} % Use the Adobe Times Roman as the default text font together with math symbols from the Sym­bol, Chancery and Com­puter Modern fonts

\usepackage{microtype} % Slightly tweak font spacing for aesthetics
\usepackage[utf8]{inputenc} % Required for including letters with accents
\usepackage[T1]{fontenc} % Use 8-bit encoding that has 256 glyphs

%----------------------------------------------------------------------------------------
%	BIBLIOGRAPHY AND INDEX
%----------------------------------------------------------------------------------------

\usepackage[style=numeric,citestyle=numeric,sorting=nyt,sortcites=true,autopunct=true,babel=hyphen,hyperref=true,abbreviate=false,backref=true,backend=biber]{biblatex}
\addbibresource{bibliography.bib} % BibTeX bibliography file
\defbibheading{bibempty}{}

\usepackage{calc} % For simpler calculation - used for spacing the index letter headings correctly
\usepackage{makeidx} % Required to make an index
\makeindex % Tells LaTeX to create the files required for indexing

%----------------------------------------------------------------------------------------
%	MAIN TABLE OF CONTENTS
%----------------------------------------------------------------------------------------

\usepackage{titletoc} % Required for manipulating the table of contents

\contentsmargin{0cm} % Removes the default margin

% Part text styling (this is mostly taken care of in the PART HEADINGS section of this file)
\titlecontents{part}
	[0cm] % Left indentation
	{\addvspace{20pt}\bfseries} % Spacing and font options for parts
	{}
	{}
	{}

% Chapter text styling
\titlecontents{chapter}
	[1.25cm] % Left indentation
	{\addvspace{12pt}\large\sffamily\bfseries} % Spacing and font options for chapters
	{\color{azul!60}\contentslabel[\Large\thecontentslabel]{1.25cm}\color{azul}} % Formatting of numbered sections of this type
	{\color{azul}} % Formatting of numberless sections of this type
	{\color{azul!60}\normalsize\;\titlerule*[.5pc]{.}\;\thecontentspage} % Formatting of the filler to the right of the heading and the page number

% Section text styling
\titlecontents{section}
	[1.25cm] % Left indentation
	{\addvspace{3pt}\sffamily\bfseries} % Spacing and font options for sections
	{\contentslabel[\thecontentslabel]{1.25cm}} % Formatting of numbered sections of this type
	{} % Formatting of numberless sections of this type
	{\hfill\color{black}\thecontentspage} % Formatting of the filler to the right of the heading and the page number

% Subsection text styling
\titlecontents{subsection}
	[1.25cm] % Left indentation
	{\addvspace{1pt}\sffamily\small} % Spacing and font options for subsections
	{\contentslabel[\thecontentslabel]{1.25cm}} % Formatting of numbered sections of this type
	{} % Formatting of numberless sections of this type
	{\ \titlerule*[.5pc]{.}\;\thecontentspage} % Formatting of the filler to the right of the heading and the page number

% Figure text styling
\titlecontents{figure}
	[1.25cm] % Left indentation
	{\addvspace{1pt}\sffamily\small} % Spacing and font options for figures
	{\thecontentslabel\hspace*{1em}} % Formatting of numbered sections of this type
	{} % Formatting of numberless sections of this type
	{\ \titlerule*[.5pc]{.}\;\thecontentspage} % Formatting of the filler to the right of the heading and the page number

% Table text styling
\titlecontents{table}
	[1.25cm] % Left indentation
	{\addvspace{1pt}\sffamily\small} % Spacing and font options for tables
	{\thecontentslabel\hspace*{1em}} % Formatting of numbered sections of this type
	{} % Formatting of numberless sections of this type
	{\ \titlerule*[.5pc]{.}\;\thecontentspage} % Formatting of the filler to the right of the heading and the page number

%----------------------------------------------------------------------------------------
%	MINI TABLE OF CONTENTS IN PART HEADS
%----------------------------------------------------------------------------------------

% Chapter text styling
\titlecontents{lchapter}
	[0em] % Left indentation
	{\addvspace{15pt}\large\sffamily\bfseries} % Spacing and font options for chapters
	{\color{azul}\contentslabel[\Large\thecontentslabel]{1.25cm}\color{azul}} % Chapter number
	{}  
	{\color{azul}\normalsize\sffamily\bfseries\;\titlerule*[.5pc]{.}\;\thecontentspage} % Page number

% Section text styling
\titlecontents{lsection}
	[0em] % Left indentation
	{\sffamily\small} % Spacing and font options for sections
	{\contentslabel[\thecontentslabel]{1.25cm}} % Section number
	{}
	{}

% Subsection text styling (note these aren't shown by default, display them by searchings this file for tocdepth and reading the commented text)
\titlecontents{lsubsection}
	[.5em] % Left indentation
	{\sffamily\footnotesize} % Spacing and font options for subsections
	{\contentslabel[\thecontentslabel]{1.25cm}}
	{}
	{}

%----------------------------------------------------------------------------------------
%	HEADERS AND FOOTERS
%----------------------------------------------------------------------------------------

\usepackage{fancyhdr} % Required for header and footer configuration

\pagestyle{fancy} % Enable the custom headers and footers

\renewcommand{\chaptermark}[1]{\markboth{\sffamily\normalsize\bfseries\chaptername\ \thechapter.\ #1}{}} % Styling for the current chapter in the header
\renewcommand{\sectionmark}[1]{\markright{\sffamily\normalsize\thesection\hspace{5pt}#1}{}} % Styling for the current section in the header

\fancyhf{} % Clear default headers and footers
\fancyhead[LE,RO]{\sffamily\normalsize\thepage} % Styling for the page number in the header
\fancyhead[LO]{\rightmark} % Print the nearest section name on the left side of odd pages
\fancyhead[RE]{\leftmark} % Print the current chapter name on the right side of even pages
%\fancyfoot[C]{\thepage} % Uncomment to include a footer

\renewcommand{\headrulewidth}{0.5pt} % Thickness of the rule under the header

\fancypagestyle{plain}{% Style for when a plain pagestyle is specified
	\fancyhead{}\renewcommand{\headrulewidth}{0pt}%
}

% Removes the header from odd empty pages at the end of chapters
\makeatletter
\renewcommand{\cleardoublepage}{
\clearpage\ifodd\c@page\else
\hbox{}
\vspace*{\fill}
\thispagestyle{empty}
\newpage
\fi}

%----------------------------------------------------------------------------------------
%	THEOREM STYLES
%----------------------------------------------------------------------------------------

\usepackage{amsmath,amsfonts,amssymb,amsthm} % For math equations, theorems, symbols, etc

\newcommand{\intoo}[2]{\mathopen{]}#1\,;#2\mathclose{[}}
\newcommand{\ud}{\mathop{\mathrm{{}d}}\mathopen{}}
\newcommand{\intff}[2]{\mathopen{[}#1\,;#2\mathclose{]}}
\renewcommand{\qedsymbol}{$\blacksquare$}
\newtheorem{notation}{Notation}[chapter]

% Boxed/framed environments
\newtheoremstyle{ocrenumbox}% Theorem style name
{0pt}% Space above
{0pt}% Space below
{\normalfont}% Body font
{}% Indent amount
{\small\bf\sffamily\color{azul}}% Theorem head font
{\;}% Punctuation after theorem head
{0.25em}% Space after theorem head
{\small\sffamily\color{azul}\thmname{#1}\nobreakspace\thmnumber{\@ifnotempty{#1}{}\@upn{#2}}% Theorem text (e.g. Theorem 2.1)
\thmnote{\nobreakspace\the\thm@notefont\sffamily\bfseries\color{black}---\nobreakspace#3.}} % Optional theorem note

\newtheoremstyle{blacknumex}% Theorem style name
{5pt}% Space above
{5pt}% Space below
{\normalfont}% Body font
{} % Indent amount
{\small\bf\sffamily}% Theorem head font
{\;}% Punctuation after theorem head
{0.25em}% Space after theorem head
{\small\sffamily{\tiny\ensuremath{\blacksquare}}\nobreakspace\thmname{#1}\nobreakspace\thmnumber{\@ifnotempty{#1}{}\@upn{#2}}% Theorem text (e.g. Theorem 2.1)
\thmnote{\nobreakspace\the\thm@notefont\sffamily\bfseries---\nobreakspace#3.}}% Optional theorem note

\newtheoremstyle{blacknumbox} % Theorem style name
{0pt}% Space above
{0pt}% Space below
{\normalfont}% Body font
{}% Indent amount
{\small\bf\sffamily}% Theorem head font
{\;}% Punctuation after theorem head
{0.25em}% Space after theorem head
{\small\sffamily\thmname{#1}\nobreakspace\thmnumber{\@ifnotempty{#1}{}\@upn{#2}}% Theorem text (e.g. Theorem 2.1)
\thmnote{\nobreakspace\the\thm@notefont\sffamily\bfseries---\nobreakspace#3.}}% Optional theorem note

% Non-boxed/non-framed environments
\newtheoremstyle{ocrenum}% Theorem style name
{5pt}% Space above
{5pt}% Space below
{\normalfont}% Body font
{}% Indent amount
{\small\bf\sffamily\color{azul}}% Theorem head font
{\;}% Punctuation after theorem head
{0.25em}% Space after theorem head
{\small\sffamily\color{azul}\thmname{#1}\nobreakspace\thmnumber{\@ifnotempty{#1}{}\@upn{#2}}% Theorem text (e.g. Theorem 2.1)
\thmnote{\nobreakspace\the\thm@notefont\sffamily\bfseries\color{black}---\nobreakspace#3.}} % Optional theorem note
\makeatother

% Defines the theorem text style for each type of theorem to one of the three styles above
\newcounter{dummy} 
\numberwithin{dummy}{section}
\theoremstyle{ocrenumbox}
\newtheorem{theoremeT}[dummy]{Theorem}
\newtheorem{problem}{Problem}[chapter]
\newtheorem{exerciseT}{Exercise}[chapter]
\theoremstyle{blacknumex}
\newtheorem{exampleT}{Example}[chapter]
\theoremstyle{blacknumbox}
\newtheorem{vocabulary}{Vocabulary}[chapter]
\newtheorem{definitionT}{Definition}[section]
\newtheorem{corollaryT}[dummy]{Corollary}
\theoremstyle{ocrenum}
\newtheorem{proposition}[dummy]{Proposition}

%----------------------------------------------------------------------------------------
%	DEFINITION OF COLORED BOXES
%----------------------------------------------------------------------------------------

\RequirePackage[framemethod=default]{mdframed} % Required for creating the theorem, definition, exercise and corollary boxes

% Theorem box
\newmdenv[skipabove=7pt,
skipbelow=7pt,
backgroundcolor=black!5,
linecolor=azul,
innerleftmargin=5pt,
innerrightmargin=5pt,
innertopmargin=5pt,
leftmargin=0cm,
rightmargin=0cm,
innerbottommargin=5pt]{tBox}

% Exercise box	  
\newmdenv[skipabove=7pt,
skipbelow=7pt,
rightline=false,
leftline=true,
topline=false,
bottomline=false,
backgroundcolor=azul!10,
linecolor=azul,
innerleftmargin=5pt,
innerrightmargin=5pt,
innertopmargin=5pt,
innerbottommargin=5pt,
leftmargin=0cm,
rightmargin=0cm,
linewidth=4pt]{eBox}	

% Definition box
\newmdenv[skipabove=7pt,
skipbelow=7pt,
rightline=false,
leftline=true,
topline=false,
bottomline=false,
linecolor=azul,
innerleftmargin=5pt,
innerrightmargin=5pt,
innertopmargin=0pt,
leftmargin=0cm,
rightmargin=0cm,
linewidth=4pt,
innerbottommargin=0pt]{dBox}	

% Corollary box
\newmdenv[skipabove=7pt,
skipbelow=7pt,
rightline=false,
leftline=true,
topline=false,
bottomline=false,
linecolor=gray,
backgroundcolor=black!5,
innerleftmargin=5pt,
innerrightmargin=5pt,
innertopmargin=5pt,
leftmargin=0cm,
rightmargin=0cm,
linewidth=4pt,
innerbottommargin=5pt]{cBox}

% Creates an environment for each type of theorem and assigns it a theorem text style from the "Theorem Styles" section above and a colored box from above
\newenvironment{theorem}{\begin{tBox}\begin{theoremeT}}{\end{theoremeT}\end{tBox}}
\newenvironment{exercise}{\begin{eBox}\begin{exerciseT}}{\hfill{\color{azul}\tiny\ensuremath{\blacksquare}}\end{exerciseT}\end{eBox}}				  
\newenvironment{definition}{\begin{dBox}\begin{definitionT}}{\end{definitionT}\end{dBox}}	
\newenvironment{example}{\begin{exampleT}}{\hfill{\tiny\ensuremath{\blacksquare}}\end{exampleT}}		
\newenvironment{corollary}{\begin{cBox}\begin{corollaryT}}{\end{corollaryT}\end{cBox}}	

%----------------------------------------------------------------------------------------
%	REMARK ENVIRONMENT
%----------------------------------------------------------------------------------------

\newenvironment{remark}{\par\vspace{10pt}\small % Vertical white space above the remark and smaller font size
\begin{list}{}{
\leftmargin=35pt % Indentation on the left
\rightmargin=25pt}\item\ignorespaces % Indentation on the right
\makebox[-2.5pt]{\begin{tikzpicture}[overlay]
\node[draw=azul!60,line width=1pt,circle,fill=azul!25,font=\sffamily\bfseries,inner sep=2pt,outer sep=0pt] at (-15pt,0pt){\textcolor{azul}{N}};\end{tikzpicture}} % Orange R in a circle
\advance\baselineskip -1pt}{\end{list}\vskip5pt} % Tighter line spacing and white space after remark

%----------------------------------------------------------------------------------------
%	SECTION NUMBERING IN THE MARGIN
%----------------------------------------------------------------------------------------

\makeatletter
\renewcommand{\@seccntformat}[1]{\llap{\textcolor{azul}{\csname the#1\endcsname}\hspace{1em}}}                    
\renewcommand{\section}{\@startsection{section}{1}{\z@}
{-4ex \@plus -1ex \@minus -.4ex}
{1ex \@plus.2ex }
{\normalfont\large\sffamily\bfseries}}
\renewcommand{\subsection}{\@startsection {subsection}{2}{\z@}
{-3ex \@plus -0.1ex \@minus -.4ex}
{0.5ex \@plus.2ex }
{\normalfont\sffamily\bfseries}}
\renewcommand{\subsubsection}{\@startsection {subsubsection}{3}{\z@}
{-2ex \@plus -0.1ex \@minus -.2ex}
{.2ex \@plus.2ex }
{\normalfont\small\sffamily\bfseries}}                        
\renewcommand\paragraph{\@startsection{paragraph}{4}{\z@}
{-2ex \@plus-.2ex \@minus .2ex}
{.1ex}
{\normalfont\small\sffamily\bfseries}}

%----------------------------------------------------------------------------------------
%	PART HEADINGS
%----------------------------------------------------------------------------------------

% Numbered part in the table of contents
\newcommand{\@mypartnumtocformat}[2]{%
	\setlength\fboxsep{0pt}%
	\noindent\colorbox{azul!20}{\strut\parbox[c][.7cm]{\ecart}{\color{azul!70}\Large\sffamily\bfseries\centering#1}}\hskip\esp\colorbox{azul!40}{\strut\parbox[c][.7cm]{\linewidth-\ecart-\esp}{\Large\sffamily\centering#2}}%
}

% Unnumbered part in the table of contents
\newcommand{\@myparttocformat}[1]{%
	\setlength\fboxsep{0pt}%
	\noindent\colorbox{azul!40}{\strut\parbox[c][.7cm]{\linewidth}{\Large\sffamily\centering#1}}%
}

\newlength\esp
\setlength\esp{4pt}
\newlength\ecart
\setlength\ecart{1.2cm-\esp}
\newcommand{\thepartimage}{}%
\newcommand{\partimage}[1]{\renewcommand{\thepartimage}{#1}}%
\def\@part[#1]#2{%
\ifnum \c@secnumdepth >-2\relax%
\refstepcounter{part}%
\addcontentsline{toc}{part}{\texorpdfstring{\protect\@mypartnumtocformat{\thepart}{#1}}{\partname~\thepart\ ---\ #1}}
\else%
\addcontentsline{toc}{part}{\texorpdfstring{\protect\@myparttocformat{#1}}{#1}}%
\fi%
\startcontents%
\markboth{}{}%
{\thispagestyle{empty}%
\begin{tikzpicture}[remember picture,overlay]%
\node at (current page.north west){\begin{tikzpicture}[remember picture,overlay]%	
\fill[azul!20](0cm,0cm) rectangle (\paperwidth,-\paperheight);
\node[anchor=north] at (4cm,-3.25cm){\color{azul!40}\fontsize{220}{100}\sffamily\bfseries\thepart}; 
\node[anchor=south east] at (\paperwidth-1cm,-\paperheight+1cm){\parbox[t][][t]{8.5cm}{
\printcontents{l}{0}{\setcounter{tocdepth}{1}}% The depth to which the Part mini table of contents displays headings; 0 for chapters only, 1 for chapters and sections and 2 for chapters, sections and subsections
}};
\node[anchor=north east] at (\paperwidth-1.5cm,-3.25cm){\parbox[t][][t]{15cm}{\strut\raggedleft\color{white}\fontsize{30}{30}\sffamily\bfseries#2}};
\end{tikzpicture}};
\end{tikzpicture}}%
\@endpart}
\def\@spart#1{%
\startcontents%
\phantomsection
{\thispagestyle{empty}%
\begin{tikzpicture}[remember picture,overlay]%
\node at (current page.north west){\begin{tikzpicture}[remember picture,overlay]%	
\fill[azul!20](0cm,0cm) rectangle (\paperwidth,-\paperheight);
\node[anchor=north east] at (\paperwidth-1.5cm,-3.25cm){\parbox[t][][t]{15cm}{\strut\raggedleft\color{white}\fontsize{30}{30}\sffamily\bfseries#1}};
\end{tikzpicture}};
\end{tikzpicture}}
\addcontentsline{toc}{part}{\texorpdfstring{%
\setlength\fboxsep{0pt}%
\noindent\protect\colorbox{azul!40}{\strut\protect\parbox[c][.7cm]{\linewidth}{\Large\sffamily\protect\centering #1\quad\mbox{}}}}{#1}}%
\@endpart}
\def\@endpart{\vfil\newpage
\if@twoside
\if@openright
\null
\thispagestyle{empty}%
\newpage
\fi
\fi
\if@tempswa
\twocolumn
\fi}

%----------------------------------------------------------------------------------------
%	CHAPTER HEADINGS
%----------------------------------------------------------------------------------------

% A switch to conditionally include a picture, implemented by Christian Hupfer
\newif\ifusechapterimage
\usechapterimagetrue
\newcommand{\thechapterimage}{}%
\newcommand{\chapterimage}[1]{\ifusechapterimage\renewcommand{\thechapterimage}{#1}\fi}%
\newcommand{\autodot}{.}
\def\@makechapterhead#1{%
{\parindent \z@ \raggedright \normalfont
\ifnum \c@secnumdepth >\m@ne
\if@mainmatter
\begin{tikzpicture}[remember picture,overlay]
\node at (current page.north west)
{\begin{tikzpicture}[remember picture,overlay]
\node[anchor=north west,inner sep=0pt] at (0,0) {\ifusechapterimage\includegraphics[width=\paperwidth]{\thechapterimage}\fi};
\draw[anchor=west] (\Gm@lmargin,-9cm) node [line width=2pt,rounded corners=15pt,draw=azul,fill=white,fill opacity=0.5,inner sep=15pt]{\strut\makebox[22cm]{}};
\draw[anchor=west] (\Gm@lmargin+.3cm,-9cm) node {\huge\sffamily\bfseries\color{black}\thechapter\autodot~#1\strut};
\end{tikzpicture}};
\end{tikzpicture}
\else
\begin{tikzpicture}[remember picture,overlay]
\node at (current page.north west)
{\begin{tikzpicture}[remember picture,overlay]
\node[anchor=north west,inner sep=0pt] at (0,0) {\ifusechapterimage\includegraphics[width=\paperwidth]{\thechapterimage}\fi};
\draw[anchor=west] (\Gm@lmargin,-9cm) node [line width=2pt,rounded corners=15pt,draw=azul,fill=white,fill opacity=0.5,inner sep=15pt]{\strut\makebox[22cm]{}};
\draw[anchor=west] (\Gm@lmargin+.3cm,-9cm) node {\huge\sffamily\bfseries\color{black}#1\strut};
\end{tikzpicture}};
\end{tikzpicture}
\fi\fi\par\vspace*{270\p@}}}

%-------------------------------------------

\def\@makeschapterhead#1{%
\begin{tikzpicture}[remember picture,overlay]
\node at (current page.north west)
{\begin{tikzpicture}[remember picture,overlay]
\node[anchor=north west,inner sep=0pt] at (0,0) {\ifusechapterimage\includegraphics[width=\paperwidth]{\thechapterimage}\fi};
\draw[anchor=west] (\Gm@lmargin,-9cm) node [line width=2pt,rounded corners=15pt,draw=azul,fill=white,fill opacity=0.5,inner sep=15pt]{\strut\makebox[22cm]{}};
\draw[anchor=west] (\Gm@lmargin+.3cm,-9cm) node {\huge\sffamily\bfseries\color{black}#1\strut};
\end{tikzpicture}};
\end{tikzpicture}
\par\vspace*{270\p@}}
\makeatother

%----------------------------------------------------------------------------------------
%	LINKS
%----------------------------------------------------------------------------------------

\usepackage{hyperref}
\hypersetup{hidelinks,backref=true,pagebackref=true,hyperindex=true,colorlinks=false,breaklinks=true,urlcolor=azul,bookmarks=true,bookmarksopen=false}

\usepackage{bookmark}
\bookmarksetup{
open,
numbered,
addtohook={%
\ifnum\bookmarkget{level}=0 % chapter
\bookmarksetup{bold}%
\fi
\ifnum\bookmarkget{level}=-1 % part
\bookmarksetup{color=azul,bold}%
\fi
}
}
 % Insert the commands.tex file which contains the majority of the structure behind the template

%\hypersetup{pdftitle={Title},pdfauthor={Author}} % Uncomment and fill out to include PDF metadata for the author and title of the book

%----------------------------------------------------------------------------------------

\begin{document}
\renewcommand{\lstlistingname}{Programa}% Listing -> Programa

%----------------------------------------------------------------------------------------
%	TITLE PAGE
%----------------------------------------------------------------------------------------

\begingroup
\thispagestyle{empty} % Suppress headers and footers on the title page
\begin{tikzpicture}[remember picture,overlay]
\node[inner sep=0pt] (background) at (current page.center) {\includegraphics[width=\paperwidth]{background.pdf}};
\draw (current page.center) node [fill=ocre!30!white,fill opacity=0.6,text opacity=1,inner sep=1cm]{\Huge\centering\bfseries\sffamily\parbox[c][][t]{\paperwidth}{\centering Ponteiros e Arrays em C\\[15pt] % Book title
{\Large Tutorial definitivo}\\[20pt] % Subtitle
{\huge por Ted Jensen\\ tradução João Araujo}}}; % Author name
\end{tikzpicture}
\vfill
\endgroup

%----------------------------------------------------------------------------------------
%	COPYRIGHT PAGE
%----------------------------------------------------------------------------------------

\newpage
~\vfill
\thispagestyle{empty}

\noindent Domínio Público \copyright\ 2021\\ % Copyright notice

\noindent \textsc{Publicado por João Araujo}\\ % Publisher

\noindent \textsc{www.araujo.eng.uerj.br}\\ % URL

\noindent Este material é colocado em domínio público. Disponível em vários formatos em\\ http://pweb.netcom.com/~tjensen/ptr/cpoint.htm\\ e em\\ http://www.araujo.eng.uerj.br \\ % License information, replace this with your own license (if any)

%\noindent \textit{First printing, March 2019} % Printing/edition date

%----------------------------------------------------------------------------------------
%	TABLE OF CONTENTS
%----------------------------------------------------------------------------------------

%\usechapterimagefalse % If you don't want to include a chapter image, use this to toggle images off - it can be enabled later with \usechapterimagetrue

\chapterimage{chapter_head_1.pdf} % Table of contents heading image

\pagestyle{empty} % Disable headers and footers for the following pages

\tableofcontents % Print the table of contents itself

\cleardoublepage % Forces the first chapter to start on an odd page so it's on the right side of the book

\pagestyle{fancy} % Enable headers and footers again

%----------------------------------------------------------------------------------------
%	PART
%----------------------------------------------------------------------------------------

%\part{Part One}

%----------------------------------------------------------------------------------------
%	CHAPTER 1
%----------------------------------------------------------------------------------------

\chapterimage{chapter_head_2.pdf} % Chapter heading image

\chapter*{Prefácio}

Este documento tem como objetivo apresentar dicas para programadores iniciantes na linguagem de programação C. Ao longo de vários anos lendo e contribuindo para várias conferências sobre C, incluindo as da \textit{FidoNet} e \textit{UseNet}, observei que um grande número de novatos em C parece ter dificuldade em compreender os fundamentos dos indicadores. Portanto, empreendi a tarefa de tentar explicá-los em linguagem simples, com muitos exemplos.

A primeira versão deste documento foi colocada em domínio público, assim como esta. Ele foi escolhido por \textit{Bob Stout}, que o incluiu como um arquivo chamado \textbf{PTR-HELP.TXT} em sua coleção amplamente distribuída de \textbf{SNIPPETS}. Desde aquele lançamento original de 1995, adicionei uma quantidade significativa de material e fiz algumas pequenas correções no trabalho original.

Posteriormente, postei uma versão HTML por volta de 1998 em meu site em:\\
http://pweb.netcom.com/~tjensen/ptr/cpoint.htm

Depois de inúmeros pedidos, finalmente saí com esta versão em PDF que é idêntica à versão em HTML citada acima e que pode ser obtida no mesmo site.

\section*{Agradecimentos}
Há tantas pessoas que, sem saber, contribuíram para este trabalho por causa das perguntas que colocaram no \textit{FidoNet C Echo}, ou no \textit{UseNet Newsgroup comp.lang.c}, ou em várias outras conferências em outras redes, que seria impossível listar todas elas. Agradecimentos especiais a \textit{Bob Stout}, que teve a gentileza de incluir a primeira versão deste material em seu arquivo \textbf{SNIPPETS}.

\section*{Sobre o Autor}
Ted Jensen é um engenheiro eletrônico aposentado que trabalhou como designer de hardware ou gerente de designers de hardware na área de gravação magnética. Programar tem sido um hobby seu, intermitentemente, desde 1968, quando ele aprendeu a perfurar cartões para enviá-los a um mainframe. (O mainframe tinha 64K de memória de núcleo magnético!).

\section*{Uso deste Material}
Tudo aqui contido é liberado para o domínio público. Qualquer pessoa pode copiar ou distribuir este material da maneira que desejar. A única coisa que peço é que se este material for usado como um auxiliar de ensino em uma aula, eu agradeceria se fosse distribuído na íntegra, ou seja, incluindo todos os capítulos, o prefácio e a introdução. Eu também agradeceria se, em tais circunstâncias, o instrutor de tal classe me deixasse uma nota em um dos endereços abaixo informando-me sobre isso. Escrevi isso com a esperança de que seja útil a outras pessoas e, como não estou pedindo nenhuma remuneração financeira, a única maneira de saber que atingi, pelo menos parcialmente, esse objetivo é por meio do feedback de quem considera este material útil.

A propósito, você não precisa ser um instrutor ou professor para entrar em contato comigo. Agradeço uma nota de qualquer pessoa que considere o material útil ou que tenha uma crítica construtiva a oferecer. Também estou disposto a responder a perguntas enviadas por e-mail nos endereços abaixo.

\begin{verbatim}
Ted Jensen
Redwood City, California
tjensen@ix.netcom.com
July 1998
\end{verbatim}
\chapterimage{edisons-electric-pen-1600.jpg} % Chapter heading image

\chapter*{Introdução}
Se você deseja ser proficiente na escrita de código na linguagem de programação C, deve ter um conhecimento prático completo de como usar ponteiros. Infelizmente, os ponteiros C parecem representar um obstáculo para os recém-chegados, particularmente aqueles vindos de outras linguagens de computador, como Python, Java ou Javascript.

Para ajudar esses recém-chegados a compreender as dicas, escrevi o seguinte material. Para obter o máximo benefício deste material, considero importante que o usuário seja capaz de executar o código nas várias listagens contidas no artigo. Tentei, portanto, manter todo o código compatível com ANSI para que funcione com qualquer compilador compatível com ANSI. Também tentei bloquear cuidadosamente o código dentro do texto. Dessa forma, com a ajuda de um editor de texto ASCII, você pode copiar um determinado bloco de código para um novo arquivo e compilá-lo em seu sistema. Recomendo que os leitores façam isso, pois ajudará na compreensão do material.


\chapterimage{champagne-cross-section-2.jpg} % Chapter heading image

\chapter{O que é um Ponteiro?}
\lstcodestyle
Uma daquelas coisas que os iniciantes em C acham difícil é o conceito de ponteiros. O objetivo deste texto é fornecer uma introdução aos ponteiros e seu uso para esses iniciantes.

Eu descobri que muitas vezes o principal motivo pelo qual os iniciantes têm problemas com ponteiros é que eles têm uma percepção fraca ou mínima das variáveis (como são usadas em C). Assim, começamos com uma discussão das variáveis C em geral.

Uma \textbf{variável} em um programa é algo com um nome, cujo valor pode variar. A maneira como o compilador e o \textit{linker} lidam com isso é que ele atribui um bloco específico de memória dentro do computador para armazenar o valor daquela variável. O tamanho desse bloco depende do intervalo no qual a variável pode variar. Por exemplo, em PCs, o tamanho de uma variável inteira é de 4 bytes e o de um inteiro longo é de 8 bytes. Em C, o tamanho de um tipo de variável, como um inteiro, não precisa ser o mesmo em todos os tipos de máquinas.

Quando declaramos uma variável, informamos ao compilador duas coisas, o nome da variável e o tipo da variável. Por exemplo, declaramos uma variável do tipo inteiro com o nome \textit{k} escrevendo:

\begin{lstlisting}
	int k;
\end{lstlisting}

Ao ver a parte ``int'' dessa instrução, o compilador reserva 4 bytes de memória (em um PC de 32 bits) para armazenar o valor do inteiro. Também configura uma tabela de símbolos. Nessa tabela, ele adiciona o símbolo \textit{k} e o endereço relativo na memória onde esses 4 bytes foram reservados.

Então, se mais tarde escrevermos:
\begin{lstlisting}
	k = 2;
\end{lstlisting}
esperamos que, em tempo de execução quando esta instrução for executada, o valor 2 seja colocado naquele local de memória reservado para o armazenamento do valor de \textbf{k}. Em C, nos referimos a uma variável como o inteiro \textbf{k} como um ``objeto''.

Em certo sentido, existem dois `valores'' associados ao objeto \textbf{k}. Um é o valor do número inteiro armazenado lá (2 no exemplo acima) e o outro é o ``valor'' da localização da memória, ou seja, o endereço de \textbf{k}. Alguns textos referem-se a esses dois valores com a nomenclatura \textbf{\textit{rvalue}} (valor à direita, pronunciado ``ar value'') e \textbf{\textit{lvalue}} (valor à esquerda, pronunciado ``el value''), respectivamente.

Em alguns linguagens, o \textit{lvalue} é o valor permitido no lado esquerdo do operador de atribuição `=' (ou seja, o endereço onde termina o resultado da avaliação do lado direito). O \textit{rvalue} é aquele que está no lado direito da instrução de atribuição, o 2 acima. \textit{Rvalues} não podem ser usados no lado esquerdo da instrução de atribuição. Assim: \textbf{2 = k}; é ilegal.

Na verdade, a definição acima de ``lvalue'' é um pouco modificada para C. De acordo com \textbf{K\&R II} (página 197): \cite{kr}.

\begin{quotation}
	``Um \textbf{objeto} é uma região nomeada de armazenamento; um \textbf{lvalue} é uma expressão que se refere a um objeto''`
\end{quotation}

No entanto, neste ponto, a definição originalmente citada acima é suficiente. À medida que nos familiarizamos com os ponteiros, entraremos em mais detalhes sobre isso.

Ok, agora considere:

\begin{lstlisting}[numbers=left, numberstyle=\tiny]
	int j, k;

	k = 2;
	j = 7; (*@\label{linha1}@*)
	k = j; (*@\label{linha2}@*)
\end{lstlisting}

Acima, o compilador interpreta o \textbf{j} na linha \ref{linha1} como o endereço da variável \textbf{j} (seu \textit{lvalue}) e cria o código para copiar o valor 7 para esse endereço. Na linha \ref{linha2}, entretanto, \textbf{j} é interpretado como seu \textit{rvalue} (uma vez que está no lado direito do operador de atribuição '='). Ou seja, aqui \textbf{j} se refere ao valor armazenado no local da memória reservado para \textbf{j}, neste caso 7. Assim, o 7 é copiado para o endereço designado pelo \textit{lvalor} de \textbf{k}.

Em todos esses exemplos, estamos usando inteiros de 4 bytes, portanto, todas as cópias de \textit{rvalues} de um local de armazenamento para o outro são feitas copiando 4 bytes. Se estivéssemos usando números inteiros longos, estaríamos copiando 8 bytes.

Agora, digamos que temos um motivo para querer uma variável projetada para conter um \textit{lvalue} (um endereço). O tamanho necessário para armazenar tal valor depende do sistema. Em computadores de mesa mais antigos com 64 K de memória total, o endereço de qualquer ponto da memória pode estar contido em 2 bytes. Os computadores com mais memória exigiriam mais bytes para armazenar um endereço. O tamanho real necessário não é muito importante, desde que tenhamos uma maneira de informar ao compilador que o que queremos armazenar é um endereço.

Essa variável é chamada de \textbf{ponteiro} (por razões que, esperamos, ficarão mais claras um pouco mais tarde). Em C, quando definimos uma variável como ponteiro, fazemos isso precedendo seu nome com um asterisco. Em C, também damos ao nosso ponteiro um tipo que, neste caso, se refere ao tipo de dados armazenados no endereço que iremos armazenar em nosso ponteiro. Por exemplo, considere a declaração da variável:

\begin{lstlisting}
	int *ptr;
\end{lstlisting}

\textbf{ptr} é o nome da nossa variável (assim como \textbf{k} era o nome da nossa variável inteira). O `*' informa ao compilador que queremos uma variável de ponteiro, ou seja, separar quantos bytes forem necessários para armazenar um endereço na memória. O \textbf{int} diz que pretendemos usar nossa variável de ponteiro para armazenar o endereço de um inteiro. Diz-se que esse ponteiro ``aponta para'' um número inteiro. No entanto, observe que quando escrevemos \textbf{int k;} não atribuímos um valor a \textbf{k}. Se essa definição for feita fora de qualquer função, os compiladores compatíveis com ANSI irão inicializá-la com zero. Da mesma forma, \textbf{ptr} não tem valor, ou seja, não armazenamos um endereço nele na declaração acima. Nesse caso, novamente se a declaração estiver fora de qualquer função, ela é inicializada com um valor garantido de forma que não aponte para nenhum objeto ou função C. Um ponteiro inicializado dessa maneira é chamado de ponteiro ``nulo''.

O padrão de bits real usado para um ponteiro nulo pode ou não ser avaliado como zero, pois depende do sistema específico no qual o código é desenvolvido. Para tornar o código-fonte compatível entre vários compiladores em vários sistemas, uma macro é usada para representar um ponteiro nulo. Essa macro é chamada de \textbf{NULL}. Portanto, definir o valor de um ponteiro usando a macro \textbf{NULL}, como em uma instrução de atribuição como \textbf{ptr = NULL}, garante que o ponteiro se tornou um ponteiro nulo. Da mesma forma, assim como se pode testar um valor inteiro igual a zero, como em \textbf{if (k == 0)}, podemos testar um ponteiro nulo usando \textbf{if (ptr == NULL)}.

Mas, de volta ao uso de nossa nova variável \textbf{ptr}. Suponha agora que queremos armazenar em \textbf{ptr} o endereço de nossa variável inteira \textbf{k}. Para fazer isso, usamos o operador \textbf{\&} unário e escrevemos:

\begin{lstlisting}
	ptr = &k;
\end{lstlisting}

O que o operador \textbf{\&} faz é recuperar o \textit{lvalue} (endereço) de \textbf{k}, embora \textbf{k} esteja no lado direito do operador de atribuição `=', e o copia para o conteúdo do nosso ponteiro \textbf{ptr}. Agora, diz-se que \textbf{ptr} ``aponta para'' \textbf{k}. Tenha paciência conosco agora, há apenas mais um operador que precisamos discutir.

O ``operador de desreferenciação'' é o asterisco e é usado da seguinte forma:

\begin{lstlisting}
	*ptr = 7;
\end{lstlisting}
irá copiar 7 para o endereço apontado por \textbf{ptr}. Assim, se \textbf{ptr} ``aponta para'' (contém o endereço de) \textbf{k}, a instrução acima irá definir o valor de \textbf{k} para 7. Ou seja, quando usamos o `*' desta forma estamos nos referindo ao valor para o qual \textbf{ptr} está apontando, não o valor do próprio ponteiro.

Da mesma forma, poderíamos escrever:
\begin{lstlisting}
	printf("%d\n",*ptr);
\end{lstlisting}
para imprimir na tela o valor inteiro armazenado no endereço apontado por \textbf{ptr;}.

Uma maneira de ver como tudo isso se encaixa seria executar o programa a seguir e, em seguida, revisar o código e a saída com cuidado.

\lstinputlisting[label={prog1-1}, caption={Variáveis e ponteiros}, numbers=left, numberstyle=\tiny, stepnumber=1]{code/prog1-1.c}

Um resultado possível, executando em um computador de 64 bits, seria:
\lstconsolestyle
\begin{lstlisting}
j tem valor 1 e está armazenado em 0x562dcb413018
k tem valor 2 e está armazenado em 0x562dcb413028
ptr tem valor 0x562dcb413028 e está armazenado em 0x562dcb413020
O valor do inteiro apontado por ptr é 2
\end{lstlisting}
\lstcodestyle
\begin{remark}
	Ainda temos que discutir aqueles aspectos de C que requerem o uso da expressão \textbf{(void *)} usada aqui. Por enquanto, inclua-o em seu código de teste. Explicaremos a razão por trás dessa expressão mais tarde.

\end{remark}

A Figura \ref{enderecos} representa esquematicamente a memória para este arranjo de variáveis. Lembre-se que neste caso estamos trabalhando com um computador de 64 bits, assim os endereços têm 8 bytes. 

\begin{figure}[ht]
	\begin{center}
		% Graphic for TeX using PGF
% Title: /home/araujo/Dropbox/VerbTeX/LivroComputadores/Pictures/memoria.dia
% Creator: Dia v0.97.3
% CreationDate: Sun May 21 11:59:17 2017
% For: araujo
% \usepackage{tikz}
% The following commands are not supported in PSTricks at present
% We define them conditionally, so when they are implemented,
% this pgf file will use them.
\ifx\du\undefined
  \newlength{\du}
\fi
\setlength{\du}{15\unitlength}
\begin{tikzpicture}
\pgftransformxscale{1.000000}
\pgftransformyscale{-1.000000}
\definecolor{dialinecolor}{rgb}{0.000000, 0.000000, 0.000000}
\pgfsetstrokecolor{dialinecolor}
\definecolor{dialinecolor}{rgb}{1.000000, 1.000000, 1.000000}
\pgfsetfillcolor{dialinecolor}

\draw (31\du,9.4\du)--(31.\du,8\du)--(25\du,8\du)--(25\du,9.4\du);
\node at (27.5\du,8.8\du){1};
\draw (31\du,10.8\du)--(31.\du,9.4\du)--(25\du,9.4\du)--(25\du,10.8\du);
\node at (28\du,10.2\du){0x562dcb413028};
\draw (31\du,12.2\du)--(31.\du,10.8\du)--(25\du,10.8\du)--(25\du,12.2\du);
\node at (27.5\du,11.6\du){2};
\draw (31\du,13.6\du)--(31.\du,12.2\du)--(25\du,12.2\du)--(25\du,13.6\du);
%\node at (27.5\du,13\du){0110 0100};
%\draw (31\du,15\du)--(31.\du,13.6\du)--(25\du,13.6\du)--(25\du,15\du);
%\node at (27.5\du,14.4\du){0110 1100};
%\draw (31\du,16.4\du)--(31.\du,15\du)--(25\du,15\du)--(25\du,16.4\du);
%\node at (27.5\du,15.8\du){0110 1111};
%\draw (31\du,17.8\du)--(31.\du,16.4\du)--(25\du,16.4\du)--(25\du,17.8\du);
%\node at (27.5\du,17.2\du){0110 1100};
%\draw (31\du,17.8\du)--(31.\du,16.4\du)--(25\du,16.4\du)--(25\du,17.8\du);
%\node at (27.5\du,18.6\du){0110 1100};
%\draw (31\du,19.2\du)--(31.\du,17.8\du)--(25\du,17.8\du)--(25\du,19.2\du)--cycle;
\node[anchor=west] at (27.5\du,12.7\du){.};
\node[anchor=west] at (27.5\du,13.100000\du){.};
\node[anchor=west] at (27.5\du,13.600000\du){.};

\definecolor{dialinecolor}{rgb}{0.000000, 0.000000, 0.000000}
\pgfsetstrokecolor{dialinecolor}
%\draw (25\du,19.2\du)--(25\du,20.6\du);
%\draw (31\du,19.2\du)--(31\du,20.6\du);
\draw (25\du,6.6\du)--(25\du,8\du);
\draw (31\du,6.6\du)--(31\du,8\du);
\node[anchor=west] at (27.5\du,7.600000\du){.};
\node[anchor=west] at (27.5\du,7.000000\du){.};
\node[anchor=west] at (27.5\du,6.600000\du){.};
\node[anchor=west] at (19.5\du,8.8\du){0x562dcb413018};
\node[anchor=west] at (19.5\du,10.2\du){0x562dcb413020};
\node[anchor=west] at (19.5\du,11.6\du){0x562dcb413028};
%\node[anchor=west] at (19.5\du,13\du){1204};
%\node[anchor=west] at (19.5\du,14.4\du){1205};
%\node[anchor=west] at (19.5\du,15.8\du){1206};
%\node[anchor=west] at (19.5\du,17.2\du){1207};
%\node[anchor=west] at (19.5\du,18.4\du){1208};
\node[anchor=west] at (19.7\du,6.002347\du) {Endere\c{c}o};
\node[anchor=west] at (26.3\du,6.002347\du){Mem\'{o}ria};
\node[anchor=west] at (31.3\du,6.002347\du){Vari\'{a}vel};
\node[anchor=west] at (31.3\du,8.8\du){j};
\node[anchor=west] at (31.3\du,10.2\du){ptr};
\node[anchor=west] at (31.3\du,11.6\du){k};
\end{tikzpicture}

		\caption{Endereços e memória.}
		\label{fig:esquemamemoria}
	\end{center}
\label{enderecos}
\end{figure}


\section*{Para revisar:}
\begin{itemize}
	\item Uma variável é declarada dando-lhe um tipo e um nome (por exemplo, \textbf{int k;})
	\item Uma variável de ponteiro é declarada dando a ela um tipo e um nome (por exemplo, \textbf{int * ptr}) onde o asterisco diz ao compilador que a variável chamada \textbf{ptr} é uma variável de ponteiro e o tipo diz ao compilador para qual tipo o ponteiro deve apontar (inteiro nesse caso).
	\item Uma vez que uma variável é declarada, podemos obter seu endereço precedendo seu nome com o operador unário \textbf{\&}, como em \textbf{\&k}.
	\item Podemos ``desreferenciar'' um ponteiro, ou seja, referir-se ao valor para o qual ele aponta, usando o operador unário `*' como em \textbf{*ptr}.
	\item Um ``\textit{lvalue}'' de uma variável é o valor do seu endereço, ou seja, onde está armazenado na memória. O ``\textit{rvalue}'' de uma variável é o valor armazenado nessa variável (naquele endereço).
\end{itemize}


\chapterimage{operator.jpg} % Chapter heading image

\chapter{Ponteiros e Arrays}
Ok, vamos em frente. Vamos considerar porque precisamos identificar o tipo de variável para a qual um ponteiro aponta, como em:
\begin{lstlisting}
	int *ptr;
\end{lstlisting}

Uma razão para fazer isso é que mais tarde, uma vez que \textbf{ptr} ``aponta para'' algo, se escrevermos:
\begin{lstlisting}
	*ptr = 2;
\end{lstlisting}
o compilador saberá quantos bytes copiar para aquele local de memória apontado por \textbf{ptr}. Se \textbf{ptr} for declarado apontando para um inteiro, 4 bytes serão copiados, se for longo, 8 bytes serão copiados. Da mesma forma, para \textit{floats} e \textit{doubles}, o número apropriado será copiado. Porém, definir o tipo para o qual o ponteiro aponta permite uma série de outras maneiras interessantes que um compilador pode interpretar o código. Por exemplo, considere um bloco na memória consistindo em dez inteiros em uma linha. Ou seja, 40 bytes de memória são reservados para conter 10 inteiros.

Agora, digamos que apontamos nosso ponteiro de inteiro \textbf{ptr} para o primeiro desses inteiros. Além disso, digamos que o inteiro esteja localizado na posição de memória 100 (decimal). O que acontece quando escrevemos:
\begin{lstlisting}
	ptr + 1;
\end{lstlisting}

Porque o compilador ``sabe'' que este é um ponteiro (ou seja, seu valor é um endereço) e que aponta para um inteiro (seu endereço atual, 100, é o endereço de um inteiro), ele adiciona 4 a \textbf{ptr} em vez de 1, então o ponteiro ``aponta para'' o \textbf{próximo inteiro}, no local da memória 104. Da mesma forma, se o \textbf{ptr} fosse declarado como um ponteiro para um longo, ele adicionaria 8 a ele em vez de 1.
O mesmo vale para outros tipos de dados, como floats, doubles ou até mesmo tipos de dados definidos pelo usuário, como estruturas. Obviamente, esse não é o mesmo tipo de ``adição'' em que normalmente pensamos. Em C, é referido como adição usando ``aritmética de ponteiros'', um termo ao qual voltaremos mais tarde.

Da mesma forma, uma vez que \textbf{++ptr} e \textbf{ptr++} são ambos equivalentes a \textbf{ptr + 1} (embora o ponto no programa quando \textbf{ptr} é incrementado possa ser diferente), incrementar um ponteiro usando o operador unário \textbf{++}, seja pré- ou pós-, incrementa o endereço que ele armazena pela quantidade \textit{sizeof(tipo)} onde ``tipo'' é o tipo do objeto apontado. (ou seja, 4 para um número inteiro, 8 para um longo, etc.).

Como um bloco de 10 inteiros localizado contiguamente na memória é, por definição, um array de inteiros, isso traz uma relação interessante entre arrays e ponteiros.

Considere o seguinte:
\begin{lstlisting}
	int meu_array[] = {1,23,17,4,-5,100};
\end{lstlisting}

Aqui temos um array contendo 6 inteiros. Referimo-nos a cada um desses inteiros por meio de um subscrito para \textbf{meu\_array}, ou seja, usando de \textbf{meu\_array[0]} até \textbf{meu\_array[5]}. Mas, podemos acessá-los alternativamente por meio de um ponteiro da seguinte maneira:

\begin{lstlisting}
	int *ptr;
	ptr = &meu_array[0]; /* o ponteiro aponta para o primeiro
	                        inteiro em nosso array */
\end{lstlisting}

E então poderíamos imprimir nosso array usando a notação de array ou desreferenciando nosso ponteiro. O código a seguir ilustra isso:

\lstinputlisting[caption=Ponteiros e arrays, numbers=left]{code/prog2-1.c}

Compile e execute o programa acima e observe cuidadosamente as linhas A e B e que o programa imprime os mesmos valores em ambos os casos. Observe também como desreferenciamos nosso ponteiro na linha B, ou seja, primeiro adicionamos \textbf{i} a ele e, em seguida, desreferenciamos o novo ponteiro.
Se você digitou o programa corretamente, sua sáida deve ser algo como:
\lstconsolestyle
\begin{lstlisting}
meu_array[0] = 1	ptr + 0 = 1
meu_array[1] = 23	ptr + 1 = 23
meu_array[2] = 17	ptr + 2 = 17
meu_array[3] = 4	ptr + 3 = 4
meu_array[4] = -5	ptr + 4 = -5
meu_array[5] = 100	ptr + 5 = 100
\end{lstlisting}
\lstcodestyle

Mude a linha B para:
\begin{lstlisting}
	printf("ptr + %d = %d\n",i, *ptr++);
\end{lstlisting}
e execute-o novamente ... depois mude para:
\begin{lstlisting}
	printf("ptr + %d = %d\n",i, *(++ptr));
\end{lstlisting}
e tente mais uma vez. A cada vez, tente prever o resultado e observe cuidadosamente o resultado real.

Em C, o padrão afirma que onde quer que possamos usar \textbf{\&var\_nome[0]}, podemos substituir isso por \textbf{var\_nome}, portanto, em nosso código onde escrevemos:
\begin{lstlisting}
	ptr = &meu_array[0];
\end{lstlisting}
podemos escrever:
\begin{lstlisting}
	ptr = meu_array;
\end{lstlisting}
para alcançar o mesmo resultado.

Isso leva muitos textos a afirmar que o nome de um array é um ponteiro. Eu prefiro pensar mentalmente ``o nome do array é o endereço do primeiro elemento no array". Muitos iniciantes (incluindo eu quando estava aprendendo) tendem a ficar confusos pensando nisso como um ponteiro. Por exemplo, enquanto podemos escrever
\begin{lstlisting}
	ptr = meu_array;
\end{lstlisting}
mas não podemos escrever
\begin{lstlisting}
	meu_array = ptr;
\end{lstlisting}

O motivo é que, embora \textbf{ptr} seja uma variável, \textbf{meu\_array} é uma constante. Ou seja, o local em que o primeiro elemento de \textbf{meu\_array} será armazenado não pode ser alterado depois que \textbf{meu\_array[]} for declarado.

A Figura \ref{fig:pontarrays} mostra uma possível configuração dessas variáveis.
\begin{figure}[ht]
	\begin{center}
		% Graphic for TeX using PGF
% Title: /home/araujo/Dropbox/VerbTeX/LivroComputadores/Pictures/memoria.dia
% Creator: Dia v0.97.3
% CreationDate: Sun May 21 11:59:17 2017
% For: araujo
% \usepackage{tikz}
% The following commands are not supported in PSTricks at present
% We define them conditionally, so when they are implemented,
% this pgf file will use them.
\ifx\du\undefined
  \newlength{\du}
\fi
\setlength{\du}{15\unitlength}
\begin{tikzpicture}
\pgftransformxscale{1.000000}
\pgftransformyscale{-1.000000}
\definecolor{dialinecolor}{rgb}{0.000000, 0.000000, 0.000000}
\pgfsetstrokecolor{dialinecolor}
\definecolor{dialinecolor}{rgb}{1.000000, 1.000000, 1.000000}
\pgfsetfillcolor{dialinecolor}

\draw (31\du,9.4\du)--(31.\du,8\du)--(25\du,8\du)--(25\du,9.4\du);
\node at (27.5\du,8.8\du){1};
\draw (31\du,10.8\du)--(31.\du,9.4\du)--(25\du,9.4\du)--(25\du,10.8\du);
\node at (27.5\du,10.2\du){23};
\draw (31\du,12.2\du)--(31.\du,10.8\du)--(25\du,10.8\du)--(25\du,12.2\du);
\node at (27.5\du,11.6\du){17};
\draw (31\du,13.6\du)--(31.\du,12.2\du)--(25\du,12.2\du)--(25\du,13.6\du);
\node at (27.5\du,13\du){4};
\draw (31\du,15\du)--(31.\du,13.6\du)--(25\du,13.6\du)--(25\du,15\du);
\node at (27.5\du,14.4\du){-5};
\draw (31\du,16.4\du)--(31.\du,15\du)--(25\du,15\du)--(25\du,16.4\du);
\node at (27.5\du,15.8\du){100};
\draw (31\du,17.8\du)--(31.\du,16.4\du)--(25\du,16.4\du)--(25\du,17.8\du);
\node at (28\du,18.4\du){0x5624ea691010};
\draw (31\du,17.8\du)--(31.\du,16.4\du)--(25\du,16.4\du)--(25\du,17.8\du);
%\node at (27.5\du,18.6\du){0110 1100};
\draw (31\du,19.2\du)--(31.\du,17.8\du)--(25\du,17.8\du)--(25\du,19.2\du)--cycle;
\node[anchor=west] at (27.5\du,20.6\du){.};
\node[anchor=west] at (27.5\du,20\du){.};
\node[anchor=west] at (27.5\du,19.6\du){.};

\definecolor{dialinecolor}{rgb}{0.000000, 0.000000, 0.000000}
\pgfsetstrokecolor{dialinecolor}
\draw (25\du,19.2\du)--(25\du,20.6\du);
\draw (31\du,19.2\du)--(31\du,20.6\du);
\draw (25\du,6.6\du)--(25\du,8\du);
\draw (31\du,6.6\du)--(31\du,8\du);
\node[anchor=west] at (27.5\du,7.600000\du){.};
\node[anchor=west] at (27.5\du,7.000000\du){.};
\node[anchor=west] at (27.5\du,6.600000\du){.};
\node[anchor=west] at (19.5\du,8.8\du){0x5624ea691010};
\node[anchor=west] at (19.5\du,10.2\du){0x5624ea691014};
\node[anchor=west] at (19.5\du,11.6\du){0x5624ea691018};
\node[anchor=west] at (19.5\du,13\du){0x5624ea69101c};
\node[anchor=west] at (19.5\du,14.4\du){0x5624ea691020};
\node[anchor=west] at (19.5\du,15.8\du){0x5624ea691024};
\node[anchor=west] at (19.5\du,17.2\du){0x5624ea691028};
\node[anchor=west] at (19.5\du,18.4\du){0x5624ea691030};
\node[anchor=west] at (19.7\du,6.002347\du) {Endere\c{c}o};
\node[anchor=west] at (26.3\du,6.002347\du){Mem\'{o}ria};
\node[anchor=west] at (31.3\du,6.002347\du){Vari\'{a}vel};
\node[anchor=west] at (31.3\du,8.8\du){meu\_array[0]};
\node[anchor=west] at (31.3\du,10.2\du){meu\_array[1]};
\node[anchor=west] at (31.3\du,11.6\du){meu\_array[2]};
\node[anchor=west] at (31.3\du,13\du){meu\_array[3]};
\node[anchor=west] at (31.3\du,14.4\du){meu\_array[4]};
\node[anchor=west] at (31.3\du,15.8\du){meu\_array[5]};
\node[anchor=west] at (31.3\du,18.4\du){ptr};
\end{tikzpicture}

		\caption{Ponteiros e arrays.}
		\label{fig:pontarrays}
	\end{center}
\end{figure}

Anteriormente, ao discutir o termo ``lvalue'', citei K\&R-2, onde afirmava:

\begin{quotation}
	``Um \textbf{objeto} é uma região nomeada de armazenamento; um \textbf{lvalue} é uma expressão que se refere a um objeto''.
\end{quotation}

Isso levanta um problema interessante. Visto que \textbf{meu\_array} é uma região nomeada de armazenamento, por que \textbf{meu\_array} na instrução de atribuição acima não é um \textit{lvalue}? Para resolver esse problema, alguns se referem a \textbf{meu\_array} como um ``\textit{lvalue} não modificável''.

Modifique o programa de exemplo anterior mudando
\begin{lstlisting}
	ptr = &meu_array[0];
\end{lstlisting}
para
\begin{lstlisting}
	ptr = meu_array;
\end{lstlisting}
e execute-o novamente para verificar se os resultados são idênticos.

Agora, vamos nos aprofundar um pouco mais na diferença entre os nomes \textbf{ptr} e \textbf{meu\_array} como usados acima. Alguns escritores se referem ao nome de um array como um ponteiro constante. O que queremos dizer com isso? Bem, para entender o termo ``\textbf{\textit{constante}}'' neste sentido, vamos voltar à nossa definição do termo ``variável''. Quando declaramos uma variável, reservamos um ponto na memória para armazenar o valor do tipo apropriado. Feito isso, o nome da variável pode ser interpretado de duas maneiras. Quando usado no lado esquerdo do operador de atribuição, o compilador o interpreta como o local da memória para o qual mover aquele valor resultante da avaliação do lado direito do operador de atribuição. Mas, quando usado no lado direito do operador de atribuição, o nome de uma variável é interpretado como significando o conteúdo armazenado naquele endereço de memória reservado para conter o valor dessa variável.

Com isso em mente, vamos agora considerar a mais simples das constantes, como em:
\begin{lstlisting}
	int i, k;
	i = 2;
\end{lstlisting}

Aqui, enquanto \textbf{i} é uma variável e então ocupa espaço na porção de dados da memória, 2 é uma constante e, como tal, em vez de reservar memória no segmento de dados, é embutido diretamente no segmento de código da memória. Ou seja, ao escrever algo como \textbf{k = i;} dizemos ao compilador para criar um código que, em tempo de execução, examinará a localização da memória \textbf{\&i} para determinar o valor a ser movido para \textbf{k}, o código criado por \textbf{i = 2;} simplesmente coloca o 2 no código e não há referência ao segmento de dados. Ou seja, \textbf{k} e \textbf{i} são objetos, mas 2 não é um objeto.

Da mesma forma, no exemplo acima, uma vez que \textbf{meu\_array} é uma constante, uma vez que o compilador estabelece onde o próprio \textit{array} deve ser armazenado, ele ``sabe'' o endereço de \textbf{meu\_array[0]} e ao ver:
\begin{lstlisting}
	ptr = meu_array;
\end{lstlisting}
ele simplesmente usa esse endereço como uma constante no segmento de código e não há referência do segmento de dados além disso.

Este pode ser um bom lugar para explicar melhor o uso da expressão \textbf{(void *)} usada no Programa \ref{prog1-1} do Capítulo 1. Como vimos, podemos ter ponteiros de vários tipos. Até agora, discutimos ponteiros para inteiros e ponteiros para caracteres. Nos próximos capítulos, aprenderemos sobre ponteiros para estruturas e até mesmo ponteiros para ponteiros.

Também aprendemos que em diferentes sistemas o tamanho de um ponteiro pode variar. Acontece que também é possível que o tamanho de um ponteiro possa variar dependendo do tipo de dados do objeto para o qual ele aponta. Assim, como acontece com inteiros nos quais você pode ter problemas ao tentar atribuir um inteiro longo a uma variável do tipo inteiro curto, você pode ter problemas ao tentar atribuir os valores de ponteiros de vários tipos a variáveis de ponteiro de outros tipos.

Para minimizar esse problema, C fornece um ponteiro do tipo \textbf{void}. Podemos declarar tal ponteiro escrevendo:
\begin{lstlisting}
	void *vptr;
\end{lstlisting}

Um ponteiro \textbf{void} é uma espécie de ponteiro genérico. Por exemplo, enquanto C não permite a comparação de um ponteiro para um tipo inteiro com um ponteiro para um tipo caractere, por exemplo, qualquer um deles pode ser comparado a um ponteiro \textbf{void}. Claro, como com outras variáveis, os \textit{casts} podem ser usados para converter de um tipo de ponteiro para outro nas circunstâncias adequadas. No Programa \ref{prog1-1} do Capítulo 1, converto os ponteiros para inteiros em ponteiros vazios para torná-los compatíveis com a especificação de conversão \textbf{\%p}. Em capítulos posteriores, outros \textit{casts} serão feitos pelas razões aqui definidas.

Bem, isso é muito material técnico para digerir e não espero que um iniciante entenda tudo isso na primeira leitura. Com o tempo e experimentação, você vai querer voltar e reler os primeiros 2 capítulos. Mas, por enquanto, vamos prosseguir para a relação entre ponteiros, arrays de caracteres e strings.



%------------------------------------------------

\section{Citation}\index{Citation}

This statement requires citation \cite{article_key}; this one is more specific \cite[162]{book_key}.

%------------------------------------------------

\section{Lists}\index{Lists}

Lists are useful to present information in a concise and/or ordered way\footnote{Footnote example...}.

\subsection{Numbered List}\index{Lists!Numbered List}

\begin{enumerate}
\item The first item
\item The second item
\item The third item
\end{enumerate}

\subsection{Bullet Points}\index{Lists!Bullet Points}

\begin{itemize}
\item The first item
\item The second item
\item The third item
\end{itemize}

\subsection{Descriptions and Definitions}\index{Lists!Descriptions and Definitions}

\begin{description}
\item[Name] Description
\item[Word] Definition
\item[Comment] Elaboration
\end{description}

%----------------------------------------------------------------------------------------
%	CHAPTER 2
%----------------------------------------------------------------------------------------

\chapter{In-text Elements}

\section{Theorems}\index{Theorems}

This is an example of theorems.

\subsection{Several equations}\index{Theorems!Several Equations}
This is a theorem consisting of several equations.

\begin{theorem}[Name of the theorem]
In $E=\mathbb{R}^n$ all norms are equivalent. It has the properties:
\begin{align}
& \big| ||\mathbf{x}|| - ||\mathbf{y}|| \big|\leq || \mathbf{x}- \mathbf{y}||\\
&  ||\sum_{i=1}^n\mathbf{x}_i||\leq \sum_{i=1}^n||\mathbf{x}_i||\quad\text{where $n$ is a finite integer}
\end{align}
\end{theorem}

\subsection{Single Line}\index{Theorems!Single Line}
This is a theorem consisting of just one line.

\begin{theorem}
A set $\mathcal{D}(G)$ in dense in $L^2(G)$, $|\cdot|_0$. 
\end{theorem}

%------------------------------------------------

\section{Definitions}\index{Definitions}

This is an example of a definition. A definition could be mathematical or it could define a concept.

\begin{definition}[Definition name]
Given a vector space $E$, a norm on $E$ is an application, denoted $||\cdot||$, $E$ in $\mathbb{R}^+=[0,+\infty[$ such that:
\begin{align}
& ||\mathbf{x}||=0\ \Rightarrow\ \mathbf{x}=\mathbf{0}\\
& ||\lambda \mathbf{x}||=|\lambda|\cdot ||\mathbf{x}||\\
& ||\mathbf{x}+\mathbf{y}||\leq ||\mathbf{x}||+||\mathbf{y}||
\end{align}
\end{definition}

%------------------------------------------------

\section{Notations}\index{Notations}

\begin{notation}
Given an open subset $G$ of $\mathbb{R}^n$, the set of functions $\varphi$ are:
\begin{enumerate}
\item Bounded support $G$;
\item Infinitely differentiable;
\end{enumerate}
a vector space is denoted by $\mathcal{D}(G)$. 
\end{notation}

%------------------------------------------------

\section{Remarks}\index{Remarks}

This is an example of a remark.

\begin{remark}
The concepts presented here are now in conventional employment in mathematics. Vector spaces are taken over the field $\mathbb{K}=\mathbb{R}$, however, established properties are easily extended to $\mathbb{K}=\mathbb{C}$.
\end{remark}

%------------------------------------------------

\section{Corollaries}\index{Corollaries}

This is an example of a corollary.

\begin{corollary}[Corollary name]
The concepts presented here are now in conventional employment in mathematics. Vector spaces are taken over the field $\mathbb{K}=\mathbb{R}$, however, established properties are easily extended to $\mathbb{K}=\mathbb{C}$.
\end{corollary}

%------------------------------------------------

\section{Propositions}\index{Propositions}

This is an example of propositions.

\subsection{Several equations}\index{Propositions!Several Equations}

\begin{proposition}[Proposition name]
It has the properties:
\begin{align}
& \big| ||\mathbf{x}|| - ||\mathbf{y}|| \big|\leq || \mathbf{x}- \mathbf{y}||\\
&  ||\sum_{i=1}^n\mathbf{x}_i||\leq \sum_{i=1}^n||\mathbf{x}_i||\quad\text{where $n$ is a finite integer}
\end{align}
\end{proposition}

\subsection{Single Line}\index{Propositions!Single Line}

\begin{proposition} 
Let $f,g\in L^2(G)$; if $\forall \varphi\in\mathcal{D}(G)$, $(f,\varphi)_0=(g,\varphi)_0$ then $f = g$. 
\end{proposition}

%------------------------------------------------

\section{Examples}\index{Examples}

This is an example of examples.

\subsection{Equation and Text}\index{Examples!Equation and Text}

\begin{example}
Let $G=\{x\in\mathbb{R}^2:|x|<3\}$ and denoted by: $x^0=(1,1)$; consider the function:
\begin{equation}
f(x)=\left\{\begin{aligned} & \mathrm{e}^{|x|} & & \text{si $|x-x^0|\leq 1/2$}\\
& 0 & & \text{si $|x-x^0|> 1/2$}\end{aligned}\right.
\end{equation}
The function $f$ has bounded support, we can take $A=\{x\in\mathbb{R}^2:|x-x^0|\leq 1/2+\epsilon\}$ for all $\epsilon\in\intoo{0}{5/2-\sqrt{2}}$.
\end{example}

\subsection{Paragraph of Text}\index{Examples!Paragraph of Text}

\begin{example}[Example name]
\lipsum[2]
\end{example}

%------------------------------------------------

\section{Exercises}\index{Exercises}

This is an example of an exercise.

\begin{exercise}
This is a good place to ask a question to test learning progress or further cement ideas into students' minds.
\end{exercise}

%------------------------------------------------

\section{Problems}\index{Problems}

\begin{problem}
What is the average airspeed velocity of an unladen swallow?
\end{problem}

%------------------------------------------------

\section{Vocabulary}\index{Vocabulary}

Define a word to improve a students' vocabulary.

\begin{vocabulary}[Word]
Definition of word.
\end{vocabulary}

%----------------------------------------------------------------------------------------
%	PART
%----------------------------------------------------------------------------------------

\part{Part Two}

%----------------------------------------------------------------------------------------
%	CHAPTER 3
%----------------------------------------------------------------------------------------

\chapterimage{chapter_head_1.pdf} % Chapter heading image

\chapter{Presenting Information}

\section{Table}\index{Table}

\begin{table}[h]
\centering
\begin{tabular}{l l l}
\toprule
\textbf{Treatments} & \textbf{Response 1} & \textbf{Response 2}\\
\midrule
Treatment 1 & 0.0003262 & 0.562 \\
Treatment 2 & 0.0015681 & 0.910 \\
Treatment 3 & 0.0009271 & 0.296 \\
\bottomrule
\end{tabular}
\caption{Table caption}
\label{tab:example} % Unique label used for referencing the table in-text
%\addcontentsline{toc}{table}{Table \ref{tab:example}} % Uncomment to add the table to the table of contents
\end{table}

Referencing Table \ref{tab:example} in-text automatically.

%------------------------------------------------

\section{Figure}\index{Figure}

\begin{figure}[h]
\centering\includegraphics[scale=0.5]{placeholder.jpg}
\caption{Figure caption}
\label{fig:placeholder} % Unique label used for referencing the figure in-text
%\addcontentsline{toc}{figure}{Figure \ref{fig:placeholder}} % Uncomment to add the figure to the table of contents
\end{figure}

Referencing Figure \ref{fig:placeholder} in-text automatically.

%----------------------------------------------------------------------------------------
%	BIBLIOGRAPHY
%----------------------------------------------------------------------------------------

\chapter*{Bibliography}
\addcontentsline{toc}{chapter}{\textcolor{ocre}{Bibliography}} % Add a Bibliography heading to the table of contents

%------------------------------------------------

\section*{Articles}
\addcontentsline{toc}{section}{Articles}
\printbibliography[heading=bibempty,type=article]

%------------------------------------------------

\section*{Books}
\addcontentsline{toc}{section}{Books}
\printbibliography[heading=bibempty,type=book]

%----------------------------------------------------------------------------------------
%	INDEX
%----------------------------------------------------------------------------------------

\cleardoublepage % Make sure the index starts on an odd (right side) page
\phantomsection
\setlength{\columnsep}{0.75cm} % Space between the 2 columns of the index
\addcontentsline{toc}{chapter}{\textcolor{ocre}{Index}} % Add an Index heading to the table of contents
\printindex % Output the index

%----------------------------------------------------------------------------------------

\end{document}
